\documentclass[11pt]{article}

\usepackage[margin=1in]{geometry}
\usepackage{amsmath}
\usepackage{hyperref}
\usepackage{graphicx}
\usepackage{booktabs}

\title{PoemForge Phase A: Domain-Relative Preference Probing with Compression, Similarity, and Human-Rated Poetry}
\author{}
\date{}

\begin{document}
\maketitle


\section{Abstract}

We test whether compression-progress signals can recover human poetic preference from text alone. The original hypothesis was that a compression-based value signal might act as a label-free aesthetic proxy: if a candidate poem improves a model’s ability to compress a held-out poetic domain, that improvement might correspond to human-perceived poetic value. Phase A does not support that strong claim. Generic literary domains fail to recover human preference and sometimes anti-align with it.

However, the experiment reveals a more precise and useful result. When the held-out domain is constructed from human preference structure, compression, TF-IDF similarity, and embedding similarity all recover meaningful preference signal. The robust finding is therefore not that compression uniquely discovers aesthetic value, but that the domain \(D\) induces the value landscape. Human-shaped contrastive domains make appraisal structure legible to multiple readout metrics; generic world-grounded domains do not.

Across 25 held-out Surprise K-fold runs, compression structural readouts have the strongest run-level mean correlation under no residual controls and under residualization against other human targets. Under stronger surface controls, the gap narrows substantially. Poem-level bootstrap intervals over item resampling include zero for compression-vs-baseline comparisons, so the data do not resolve a unique compression advantage at \(n=36\). The strongest defensible claim is domain-relative supervised probing: human-labeled domains recover human appraisal structure across several readout families, while generic compression does not provide a label-free aesthetic oracle.

\section{1. Introduction}

A natural way to formalize value is as expected improvement in a model’s ability to predict or compress future observations. In a creative setting, this suggests a seductive hypothesis: a poem may be valuable if it helps compress a relevant poetic world. This project began from that hypothesis.

Let a candidate artifact be \(a\), a history be \(H\), possible future observations be \(O\_i\), and a held-out domain be \(D\). The value functional can be written as:

\[
V(a \mid H) =
\sum_i q_i
\left[
L(D \mid O_i \oplus H)
-
L(D \mid O_i \oplus H \oplus a)
\right]
-
\mathrm{cost}(a).
\]

The key invariant is that the score is not the likelihood of the candidate itself. The score is the candidate’s effect on a held-out domain \(D\). This distinction matters. If \(D\) is generic literary text, the system asks whether the candidate improves compression of a broad literary distribution. If \(D\) is constructed from high-rated versus low-rated poems, the system asks whether the candidate is closer to the human preference contrast encoded in that domain.

Phase A shows that this distinction is decisive. The original label-free version of the hypothesis fails. Generic Gutenberg-derived domains do not recover human poetic preference. But preference-shaped domains recover appraisal structure across compression, TF-IDF, and embedding readouts. The value signal lives less in compression as such than in the construction of \(D\).

This reframes PoemForge. The project is not evidence for an unsupervised aesthetic oracle. It is evidence for a domain-relative probing framework in which human-labeled contrastive domains induce measurable value landscapes.

\section{2. Experimental Setup}

\subsection{2.1 Items and human targets}

The evaluation set contains human-rated poems from the Chaudhuri-style poem evaluation substrate. Each poem has target ratings for dimensions including Aesthetic Appeal, Clarity, Creativity, Felt Valence, and Surprise. The primary target in the final controlled experiments is Surprise, because it is central to the original compression-progress hypothesis and because it provides a test of whether the readout captures something beyond simple predictability.

The normalized data stage freezes item text, target ratings, surface features, and matched-control pools. These artifacts define the item universe for the rest of the pipeline.

\subsection{2.2 Domains}

We compare two broad domain types.

First, generic literary domains are sampled from Gutenberg-derived poetry-like text. These domains test the strong label-free hypothesis: if compression-progress over generic literary text is enough to recover human preference, then generic \(D\) should correlate positively with human ratings.

Second, preference-shaped domains are constructed from high-rated and low-rated poems under a held-out K-fold protocol. For a target such as Surprise, high and low preference pools are rebuilt inside each training fold, and test items are scored against the resulting contrast. This prevents direct leakage from test items into the domain construction.

\subsection{2.3 Readout metrics}

We compare three readout families:

\begin{enumerate}
\item compression or language-model readouts;
\item TF-IDF similarity readouts;
\item embedding similarity readouts.
\end{enumerate}

For preference-shaped domains, each readout produces raw, control, and structural preference scores. The structural score subtracts a matched-control effect:

\[
v_{\mathrm{struct}} = v_{\mathrm{raw}} - v_{\mathrm{ctrl}}.
\]

For contrastive preference scoring, the main quantity is:

\[
v_{\mathrm{pref\_struct}}
=
(v_{\mathrm{high\_raw}} - v_{\mathrm{low\_raw}})
-
(v_{\mathrm{high\_ctrl}} - v_{\mathrm{low\_ctrl}}).
\]

The same conceptual structure is applied to TF-IDF and embedding baselines, producing comparable structural preference readouts.

\subsection{2.4 Controls}

We report three levels of control in the regenerated run-level summaries:

\begin{enumerate}
\item no residual controls;
\item residualization against other human target dimensions;
\item residualization against other human targets plus surface features.
\end{enumerate}

The frozen Phase A artifacts also include stacked controls using other human targets, surface features, and item-level language-model predictability. These are promoted into the bootstrap uncertainty table, because the item-level bootstrap was run over the canonical Phase A fully controlled outputs.

\section{3. Results}

\subsection{3.1 Generic domains do not recover human preference}

Generic Gutenberg-derived domains do not support the label-free compression-aesthetics hypothesis. In the generated generic-domain summary, raw compression scores are negative for both accessible and formal Gutenberg variants. Structural controls do not produce a stable positive preference signal. In several conditions, the correlations remain negative or near zero.

This result is important because it blocks the strongest interpretation of the original idea. Compression-progress against arbitrary literary text is not enough to recover human poetic appraisal. The domain must be shaped in relation to the target preference structure.

See Table 3, generated as \texttt{results/tables/table\_3\_generic\_d\_summary.*}, for the generic-domain summary.

\subsection{3.2 Preference-shaped domains recover appraisal structure}

Preference-shaped domains behave differently. In held-out Surprise K-fold runs, all three readout families recover positive structural signal. Compression has the largest mean run-level structural correlation without residual controls and after residualization against other human targets.

The main readout convergence table shows:

* without residual controls, compression exceeds both embedding and TF-IDF;
* after controlling for other human targets, compression remains ahead by mean;
* after adding surface controls, the gap narrows substantially.

This pattern supports the domain-relative interpretation. Human-shaped domains induce a preference landscape that multiple metrics can read. Compression is competitive and often strongest by point estimate, but the signal is not compression-exclusive.

See Table 1, generated as \texttt{results/tables/table\_1\_readout\_convergence.*}, and Figure 1, generated as \texttt{results/figures/figure\_1\_readout\_convergence.*}.

\subsection{3.3 Item-level bootstrap does not resolve a unique compression advantage}

Run-level means lean toward compression in several partially controlled settings. But item-level bootstrap uncertainty is wide. For the key compression-vs-baseline comparisons, the observed differences are positive, but all 95\% bootstrap intervals include zero.

This means the correct interpretation is not “compression is equal to the baselines.” The correct interpretation is underresolution at the item level. With \(n=36\) poems, the data do not support a strong claim that compression uniquely outperforms TF-IDF or embedding similarity.

The same caution applies to the matched-other diagnostic. Point estimates are positive, but bootstrap intervals include zero. The matched-other mechanism remains a plausible compression-specific thread, not a resolved result.

See Table 2, generated as \texttt{results/tables/table\_2\_bootstrap\_uncertainty.*}, and Figure 2, generated as \texttt{results/figures/figure\_2\_bootstrap\_uncertainty.*}.

\subsection{3.4 Higher-rated poems are often more predictable}

The item-level NLL analysis shows that higher-rated poems tend to be more language-model predictable. Correlations between unconditional item NLL and human targets are generally negative, especially for Aesthetic Appeal, Creativity, and Felt Valence. Surprise is also negatively correlated with NLL for the smallest model and weaker for larger models.

This is a useful boundary condition. Human-rated Surprise is not equivalent to language-model confusion. In this dataset, better or more surprising poems are not simply harder for a language model to predict. This weakens naive novelty-only interpretations and reinforces the need for domain-relative contrastive structure.

See Table 4, generated as \texttt{results/tables/table\_4\_item\_nll\_correlations.*}.

\section{4. Discussion}

\subsection{4.1 What failed}

The label-free compression oracle failed. Generic literary domains do not recover human preference. This failure is not a nuisance result; it defines the boundary of the framework. The held-out domain \(D\) is not an implementation detail. It is the source of the appraisal geometry.

A compression score can only be interpreted relative to the domain being compressed. If \(D\) is generic, the score reflects generic literary predictability. If \(D\) is preference-shaped, the score reflects the contrast encoded by that preference-shaped domain.

\subsection{4.2 What survived}

The domain-relative probing framework survived. When \(D\) is built from human preference structure, multiple readouts recover appraisal signal. This includes compression, TF-IDF, and embeddings.

The convergence across readout families is a positive result. It suggests that the preference-shaped domain contains robust structure rather than a metric-specific artifact. The system is not discovering aesthetic value from nowhere. It is reading out structure induced by a supervised domain.

\subsection{4.3 Compression’s remaining role}

Compression remains interesting, but the claim must be narrower. Compression has strong point estimates in several settings, and structural compression readouts outperform baselines by mean under no controls and under other-human-target residualization. However, the poem-level bootstrap does not resolve a unique compression advantage at the item level.

The most promising compression-specific thread is the matched-other diagnostic: comparing structural matched-other scoring against explicit normalization variants. The observed differences are positive, but unresolved under bootstrap. This should be treated as a future mechanism question rather than a present conclusion.

\subsection{4.4 Implications for PoemForge}

The practical implication is that PoemForge should not be framed as a label-free aesthetic evaluator. A more accurate framing is a domain construction and readout system. Its core capability is to build, audit, and compare domains that make a creative preference legible to different metrics.

Under this framing, the valuable object is not a universal compression score. The valuable object is a reproducible preference-probing pipeline:

\begin{enumerate}
\item define a target domain;
\item construct contrastive high/low pools;
\item score candidates through multiple readouts;
\item residualize against plausible confounds;
\item quantify uncertainty at the item level;
\item report where signal is robust, unresolved, or absent.
\end{enumerate}

\section{5. Limitations}

The primary limitation is sample size. The final item-level analysis has \(n=36\) poems, which is too small to resolve fine-grained differences between readout families. The bootstrap intervals make this limitation explicit.

A second limitation is domain specificity. The positive results depend on human-shaped contrastive domains. This is appropriate for supervised probing, but it means the results should not be generalized to unsupervised aesthetic discovery.

A third limitation is that the current pipeline freezes and promotes several Phase A inferential artifacts rather than fully recomputing every expensive analysis from scratch. This is acceptable for a reproducibility bridge, but a final archival version should progressively replace frozen-result promotion with direct regeneration.

\section{6. Conclusion}

PoemForge Phase A began with a strong hypothesis: compression-progress over a held-out literary domain might provide a label-free proxy for poetic value. The evidence does not support that claim. Generic domains fail.

The surviving result is sharper. Human-labeled domains induce value landscapes. When \(D\) is constructed from human preference structure, compression, TF-IDF, and embedding readouts all recover appraisal signal. Compression is competitive and often strongest by point estimate, but item-level bootstrap uncertainty does not establish a unique compression advantage at \(n=36\).

The final claim is therefore domain-relative supervised probing, not label-free compression aesthetics. The dragon was not in the metric. It was in \(D\).

\section{Reproducibility}

All reported manuscript tables and figures are generated by the reproducibility pipeline. The pipeline freezes normalized data, domain artifacts, score artifacts, control features, Phase A inferential results, bootstrap summaries, manuscript tables, manuscript figures, and scaffold numeric diffs.

The main manuscript-facing artifacts are:

\begin{itemize}
\item \texttt{results/tables/table\_1\_readout\_convergence.*};
\item \texttt{results/tables/table\_2\_bootstrap\_uncertainty.*};
\item \texttt{results/tables/table\_3\_generic\_d\_summary.*};
\item \texttt{results/tables/table\_4\_item\_nll\_correlations.*};
\item \texttt{results/figures/figure\_1\_readout\_convergence.*};
\item \texttt{results/figures/figure\_2\_bootstrap\_uncertainty.*};
\item \texttt{results/diffs/scaffold\_numeric\_diff.csv};
\item \texttt{results/diffs/scaffold\_diff\_report.md}.
\end{itemize}

The final scaffold check verifies 39 canonical manuscript values against expected values with tolerance \(10^{-9}\). All checks pass.

\end{document}
